\documentclass[aspectratio=169, 10pt]{beamer}

% --- Pacotes de Idioma e Codificação ---
\usepackage[utf8]{inputenc}
\usepackage[T1]{fontenc}
\usepackage[brazil]{babel}
\usepackage{booktabs}
\usepackage{graphicx}
\usepackage{tikz}
\usepackage{hyperref}
\usetikzlibrary{shapes.geometric, arrows.meta, positioning}

% --- Configuração de Tema e Cores ---
\usetheme{Madrid}
\usecolortheme{whale}

% Cores customizadas modernas
\definecolor{primaryNavy}{RGB}{24, 43, 73}
\definecolor{accentTeal}{RGB}{0, 150, 136}
\definecolor{lightGray}{RGB}{245, 247, 250}

\setbeamercolor{palette primary}{bg=primaryNavy, fg=white}
\setbeamercolor{palette secondary}{bg=accentTeal, fg=white}
\setbeamercolor{palette tertiary}{bg=primaryNavy!90, fg=white}
\setbeamercolor{structure}{fg=primaryNavy}
\setbeamercolor{titlelike}{parent=palette primary}
\setbeamercolor{block title}{bg=primaryNavy, fg=white}
\setbeamercolor{block body}{bg=lightGray, fg=black}
\setbeamercolor{block title alerted}{bg=red!80!black, fg=white}
\setbeamercolor{block body alerted}{bg=red!5, fg=black}

% Remover símbolos de navegação inferiores
\setbeamertemplate{navigation symbols}{}

% --- Metadados da Apresentação ---
\title[AgendaFácil]{\textbf{AgendaFácil: Sistema Inteligente de Agendamento}}
\subtitle{Automação com No Code, Low Code e Vibecode}

\author[Equipe AgendaFácil]{
    \textbf{Discentes:} \\
    Gabriella Dias \and Gustavo Negreiros \and João Pedro Alves \and Maria Cecília
}

\institute[UDF]{
    \textbf{Centro Universitário do Distrito Federal (UDF)} \\
    Tendências em Ciência da Computação \\
    \textbf{Docente:} Profa. Kadidja Valéria
}

\date{\today}

% --- Início do Documento ---
\begin{document}

% -----------------------------------------------------------------------------
% Slide 1: Capa
% -----------------------------------------------------------------------------
\begin{frame}[plain]
    \titlepage
\end{frame}

% -----------------------------------------------------------------------------
% Slide 2: Contexto e O Desafio
% -----------------------------------------------------------------------------
\begin{frame}{Contextualização: O Desafio dos Agendamentos}
    \begin{columns}[T]
        \begin{column}{0.48\textwidth}
            \begin{block}{Cenário Tradicional}
                \begin{itemize}
                    \item Agendamentos manuais (WhatsApp, mensagens soltas, anotações).
                    \item Falhas de comunicação e esquecimento de horários.
                    \item Retrabalho para confirmar compromissos.
                \end{itemize}
            \end{block}
        \end{column}
        
        \begin{column}{0.48\textwidth}
            \begin{alertblock}{O Obstáculo Técnico}
                \begin{itemize}
                    \item Desenvolver backend tradicional (APIs, banco SQL, servidores) é \textbf{lento e custoso} para validação inicial de MVP.
                    \item Pequenas operações precisam de soluções \textbf{ágeis e de custo mínimo}.
                \end{itemize}
            \end{alertblock}
        \end{column}
    \end{columns}
\end{frame}

% -----------------------------------------------------------------------------
% Slide 3: A Proposta: AgendaFácil
% -----------------------------------------------------------------------------
\begin{frame}{A Solução: Projeto AgendaFácil}
    \begin{block}{\textbf{Objetivo Geral}}
        Construir um protótipo funcional de agendamento com confirmação automática personalizada por e-mail, utilizando plataformas visuais e inteligência artificial generativa.
    \end{block}

    \vspace{0.3cm}
    \textbf{Princípios de Desenvolvimento do Projeto:}
    \begin{itemize}
        \item \textbf{Agilidade:} Redução drástica do ciclo de desenvolvimento (\textit{Time-to-Market}).
        \item \textbf{Custo Zero:} Utilização estrita de camadas gratuitas (\textit{Free Tiers}).
        \item \textbf{Humanização:} Mensagens contextualizadas em vez de modelos estáticos e frios.
    \end{itemize}
\end{frame}

% -----------------------------------------------------------------------------
% Slide 4: O Tripé Metodológico
% -----------------------------------------------------------------------------
\begin{frame}{O Tripé Metodológico da Disciplina}
    \begin{columns}[T]
        \begin{column}{0.32\textwidth}
            \begin{block}{\textbf{1. No Code}}
                \textbf{Google Forms \& Sheets}
                \begin{itemize}
                    \item Interface visual de entrada.
                    \item Banco de dados tabular.
                    \item Zero linhas de código.
                \end{itemize}
            \end{block}
        \end{column}

        \begin{column}{0.32\textwidth}
            \begin{block}{\textbf{2. Low Code}}
                \textbf{Make (Integromat)}
                \begin{itemize}
                    \item Orquestração do fluxo.
                    \item Conectores e gatilhos lógicos.
                    \item Configuração visual de dados.
                \end{itemize}
            \end{block}
        \end{column}

        \begin{column}{0.32\textwidth}
            \begin{block}{\textbf{3. Vibecode}}
                \textbf{Gemini / LLM}
                \begin{itemize}
                    \item Diálogo iterativo por prompts.
                    \item Adaptação semântica do e-mail.
                    \item Redação humanizada.
                \end{itemize}
            \end{block}
        \end{column}
    \end{columns}
\end{frame}

% -----------------------------------------------------------------------------
% Slide 5: Camada No Code na Prática
% -----------------------------------------------------------------------------
\begin{frame}{1. Camada No Code: Interface e Persistência}
    \begin{columns}[T]
        \begin{column}{0.5\textwidth}
            \textbf{Google Forms (Frontend):}
            \begin{itemize}
                \item Coleta dados essenciais: Nome, E-mail, Data, Horário e Motivo.
                \item Interface responsiva para mobile e desktop sem esforço de estilização (CSS).
            \end{itemize}
            \vspace{0.2cm}
            \textbf{Google Sheets (Database):}
            \begin{itemize}
                \item Cada submissão vira um registro em linha.
                \item Painel acessível para o administrador visualizar as solicitações em tempo real.
            \end{itemize}
        \end{column}
        
        \begin{column}{0.46\textwidth}
            \begin{block}{Benefício Chave}
                Entrega uma experiência familiar e sem barreiras de aprendizado para o usuário final e para o administrador.
            \end{block}
        \end{column}
    \end{columns}
\end{frame}

% -----------------------------------------------------------------------------
% Slide 6: Camada Low Code na Prática
% -----------------------------------------------------------------------------
\begin{frame}{2. Camada Low Code: Orquestração no Make}
    \begin{itemize}
        \item \textbf{Gatilho (Trigger):} Monitora novas linhas na planilha do Google Sheets.
        \item \textbf{Mapeamento de Variáveis:} Captura os dados do solicitante em payload JSON estruturado.
        \item \textbf{Encadeamento:} Aciona o módulo de Inteligência Artificial e redireciona a resposta gerada para o conector do Gmail.
    \end{itemize}

    \vspace{0.4cm}
    \begin{center}
        \begin{tikzpicture}[node distance=2.2cm, auto, >=Stealth]
            \tikzstyle{box} = [rectangle, draw=primaryNavy, fill=white, thick, text centered, rounded corners, minimum height=1cm, minimum width=2.5cm]
            \node[box] (n1) {\textbf{Sheets Trigger}};
            \node[box, right of=n1, xshift=1.2cm] (n2) {\textbf{Gemini Prompt}};
            \node[box, right of=n2, xshift=1.2cm] (n3) {\textbf{Gmail Send}};
            
            \draw[->, thick, color=accentTeal] (n1) -- (n2);
            \draw[->, thick, color=accentTeal] (n2) -- (n3);
        \end{tikzpicture}
    \end{center}
\end{frame}

% -----------------------------------------------------------------------------
% Slide 7: Camada Vibecode na Prática
% -----------------------------------------------------------------------------
\begin{frame}{3. Camada Vibecode: IA e Engenharia de Prompt}
    \textbf{Como o Vibecode atuou na solução?}
    \begin{itemize}
        \item Desenvolvimento orientado por linguagem natural (sem codificação estrita de regras).
        \item O modelo de IA recebe os dados variáveis e redige uma mensagem personalizada.
    \end{itemize}

    \vspace{0.2cm}
    \begin{exampleblock}{Estrutura Básica do Prompt Empregado}
        \textit{"Você é o assistente virtual do AgendaFácil. Escreva um e-mail empático e profissional para \{\{Nome\}\}, confirmando a solicitação de atendimento no dia \{\{Data\}\} às \{\{Horário\}\}. Inclua instruções para chegar com 10 minutos de antecedência."}
    \end{exampleblock}

    \vspace{0.2cm}
    \textbf{Resultado:} Um e-mail acolhedor, sem a rigidez fria de templates convencionais.
\end{frame}

% -----------------------------------------------------------------------------
% Slide 8: Arquitetura Ponta a Ponta
% -----------------------------------------------------------------------------
\begin{frame}{Arquitetura Completa da Solução}
    \begin{center}
        \begin{tikzpicture}[node distance=1.6cm, auto, >=Stealth, font=\small]
            \tikzstyle{step} = [rectangle, draw=primaryNavy, fill=primaryNavy!10, thick, text centered, rounded corners, minimum height=1.1cm, text width=2.2cm]
            
            \node[step] (f1) {\textbf{1. Solicitante}\\Preenche Forms};
            \node[step, right of=f1, xshift=1.2cm] (f2) {\textbf{2. Sheets}\\Registra Dados};
            \node[step, right of=f2, xshift=1.2cm] (f3) {\textbf{3. Make}\\Dispara Fluxo};
            \node[step, right of=f3, xshift=1.2cm] (f4) {\textbf{4. IA / LLM}\\Gera Redação};
            \node[step, right of=f4, xshift=1.2cm] (f5) {\textbf{5. Gmail}\\Entrega E-mail};

            \draw[->, very thick, color=accentTeal] (f1) -- (f2);
            \draw[->, very thick, color=accentTeal] (f2) -- (f3);
            \draw[->, very thick, color=accentTeal] (f3) -- (f4);
            \draw[->, very thick, color=accentTeal] (f4) -- (f5);
        \end{tikzpicture}
    \end{center}
    
    \vspace{0.5cm}
    \begin{alertblock}{Validação de Segurança}
        O fluxo foi testado com e-mails reais e respostas sintéticas para assegurar que a IA não sofra alucinações e mantenha os dados de data e horário rigorosamente exatos.
    \end{alertblock}
\end{frame}

% -----------------------------------------------------------------------------
% Slide 9: Análise Crítica: 3 Vantagens
% -----------------------------------------------------------------------------
\begin{frame}{Reflexão Crítica: Vantagens Identificadas}
    \begin{enumerate}
        \item \textbf{Agilidade Extrema de Desenvolvimento (\textit{Time-to-Market}):}
        \begin{itemize}
            \item Construção e validação do protótipo funcional em menos de 48 horas.
            \item Eliminação de infraestrutura de servidor, containers ou deploy manual.
        \end{itemize}
        \vspace{0.2cm}
        \item \textbf{Custo Operacional e de Infraestrutura Zero:}
        \begin{itemize}
            \item Utilização integral de camadas gratuitas dos serviços (Google e Make).
            \item Ideal para MVPs e microempresas testarem demandas reais.
        \end{itemize}
        \vspace{0.2cm}
        \item \textbf{Hiperpersonalização Humanizada (Vibecode):}
        \begin{itemize}
            \item Comunicação muito mais empática e adaptada ao contexto informado pelo cliente.
        \end{itemize}
    \end{enumerate}
\end{frame}

% -----------------------------------------------------------------------------
% Slide 10: Análise Crítica: 3 Limitações
% -----------------------------------------------------------------------------
\begin{frame}{Reflexão Crítica: Limitações e Desafios}
    \begin{enumerate}
        \item \textbf{Ausência de Validação de Conflitos em Tempo Real:}
        \begin{itemize}
            \item O Google Forms não consulta a planilha antes do envio; dois usuários podem submeter o mesmo horário.
            \item \textit{Mitigação:} O sistema atua como solicitação prévia, cabendo ao administrador validar horários disputados.
        \end{itemize}
        \vspace{0.2cm}
        \item \textbf{Limites de Quota e Latência (\textit{Free Tier}):}
        \begin{itemize}
            \item O plano gratuito do Make possui teto de operações mensais e intervalos de verificação periódica.
        \end{itemize}
        \vspace{0.2cm}
        \item \textbf{Dependência de Terceiros e Governança (\textit{Vendor Lock-in}):}
        \begin{itemize}
            \item Dados transitam por serviços de nuvem de múltiplos provedores sem controle de infraestrutura local.
        \end{itemize}
    \end{enumerate}
\end{frame}

% -----------------------------------------------------------------------------
% Slide 11: Divisão da Equipe e Colaboração
% -----------------------------------------------------------------------------
\begin{frame}{Trabalho em Equipe e Divisão de Papéis}
    \begin{table}[h]
        \centering
        \small
        \begin{tabular}{ll}
            \toprule
            \textbf{Integrante} & \textbf{Papel e Responsabilidade Principal} \\
            \midrule
            \textbf{Maria Cecília} & Criação do formulário e estruturação da planilha (No Code) \\
            \textbf{Gabriella Dias} & Configuração da lógica e orquestração de fluxo no Make (Low Code) \\
            \textbf{João Pedro Alves} & Integração com IA e engenharia de prompts (Vibecode) \\
            \textbf{Gustavo Negreiros} & Documentação no GitHub, testes e elaboração dos slides Beamer \\
            \bottomrule
        \end{tabular}
    \end{table}
\end{frame}

% -----------------------------------------------------------------------------
% Slide 12: Próximos Passos e Conclusão
% -----------------------------------------------------------------------------
\begin{frame}{Próximos Passos e Conclusão}
    \begin{columns}[T]
        \begin{column}{0.48\textwidth}
            \textbf{Evoluções Futuras do Protótipo:}
            \begin{itemize}
                \item Integração direta com a \textbf{Google Calendar API} para bloqueio automático de horário.
                \item Notificações instantâneas via \textbf{WhatsApp API}.
                \item Painel No Code no \textbf{AppSheet} para aprovação pelo administrador com 1 clique.
            \end{itemize}
        \end{column}

        \begin{column}{0.48\textwidth}
            \begin{block}{\textbf{Conclusão}}
                O projeto demonstrou que a união de No Code, Low Code e Vibecode democratiza a inovação, permitindo que soluções reais e sofisticadas sejam criadas com agilidade ímpar.
            \end{block}
        \end{column}
    \end{columns}

    \vspace{0.5cm}
    \begin{center}
        \large \textbf{Obrigado!} \\
        \normalsize Dúvidas e Perguntas
    \end{center}
\end{frame}

\end{document}
